\documentclass[aspectratio=169]{beamer}
\usepackage[utf8]{inputenc}
\usepackage[T1]{fontenc}
\usepackage[portuguese]{babel}
\usepackage{lmodern}
\usepackage{booktabs}
\usepackage{xcolor}
\usetheme{DCC}
\graphicspath{{imgs/}}

% Numeros da execucao (competencia 2026-09, dez particoes).
% Fonte unica: data/processed/evidencia.json -- atualizar aqui e recompilar.
\newcommand{\lidos}{73.366.147}
\newcommand{\naba}{3.616.754}
\newcommand{\recorte}{15.095}
\newcommand{\leads}{7.042}

\author{Isnan Ferreira da Mota Santos}
\title{Priorização de leads de tecnologia na Bahia}
\institute{PGCOMP/UFBA --- Sistemas Inteligentes Orientados a Dados}
\date{2026.2}

\begin{document}

\begin{frame}[plain,noframenumbering]
    \titlepage
    \begin{center}
        \small Atividade 01: o que é a base. \quad Atividade 02: o que eu fiz com ela.
    \end{center}
\end{frame}

\begin{frame}{O problema e a entrega}
    \begin{itemize}
        \item \textbf{Quem usa:} o analista que escolhe quais empresas procurar.
        \item \textbf{Hoje:} filtra planilha à mão, sem critério registrado.
        \item \textbf{Entrega:} lista ordenada, com contato e motivo da posição.
    \end{itemize}

    \begin{block}{Para montar essa lista, preciso do cadastro inteiro}
        Dados abertos de CNPJ da Receita Federal --- público, atualizado uma
        vez por mês, baixado por arquivo, sem consulta em massa por API.
    \end{block}
\end{frame}

\section{Atividade 01 --- os dados}

\begin{frame}{Atividade 01 --- as cinco tabelas usadas}
    \begin{table}
        \footnotesize
        \begin{tabular}{lrrl}
            \toprule
            Tabela & Linhas & Tamanho & O que ela fornece \\
            \midrule
            Estabelecimentos 0..9 & 73.366.147 & 5,01 GiB & atividade, município, situação, contato \\
            Empresas 0..9         & 70.085.592 & 1,28 GiB & razão social e porte \\
            Simples               & 50.396.768 & 294 MiB  & marcação de MEI \\
            Cnaes                 &      1.359 & $<$1 MiB & nome da atividade \\
            Municipios            &      5.572 & $<$1 MiB & nome do município \\
            \bottomrule
        \end{tabular}
    \end{table}

    \begin{itemize}
        \item \textbf{6,57 GiB} baixados. O servidor limita cada conexão a
              ~1,1 MiB/s; com quatro em paralelo, 3,2 MiB/s.
        \item Uma empresa tem vários estabelecimentos: por isso 73,4 milhões de
              estabelecimentos para 70,1 milhões de empresas.
    \end{itemize}
\end{frame}

\begin{frame}{Atividade 01 --- do cadastro nacional até a lista}
    \begin{table}
        \small
        \begin{tabular}{llr}
            \toprule
            Etapa & Filtro aplicado & Estabelecimentos \\
            \midrule
            Cadastro nacional      & ---                                 & \lidos \\
            Recorte geográfico     & \texttt{UF = BA}                    & \naba \\
            Recorte de atividade   & CNAE nas divisões 62 e 63           & \recorte \\
            Validação              & cinco regras --- nenhum descarte    & \recorte \\
            Lista final            & situação cadastral ativa            & \textbf{\leads} \\
            \bottomrule
        \end{tabular}
    \end{table}

    \begin{itemize}
        \item Os dois recortes tiram 99,98\% da base. O que sobra é o universo
              de trabalho do analista.
        \item O último filtro é o que mais importa --- e é o assunto de dois
              slides à frente.
    \end{itemize}
\end{frame}

\begin{frame}{Atividade 01 --- do que é feito o recorte}
    \begin{columns}[T]
        \begin{column}{0.60\textwidth}
            \begin{table}
                \footnotesize
                \begin{tabular}{llr}
                    \toprule
                    CNAE & Atividade & Empresas \\
                    \midrule
                    6209 & suporte técnico e manutenção & 3.604 \\
                    6201 & desenvolvimento sob encomenda & 2.511 \\
                    6204 & consultoria em TI & 2.278 \\
                    6319 & portais e provedores de conteúdo & 1.688 \\
                    6399 & outros serviços de informação & 1.481 \\
                    6311 & tratamento de dados e hospedagem & 1.348 \\
                    6202 & software customizável & 1.168 \\
                    6203 & software não customizável & 780 \\
                    6391 & agências de notícias & 237 \\
                    \midrule
                    & Total & \recorte \\
                    \bottomrule
                \end{tabular}
            \end{table}
        \end{column}
        \begin{column}{0.36\textwidth}
            \begin{itemize}
                \item Divisão 62, serviços de TI: \textbf{10.341} --- 68\%.
                \item Divisão 63, serviços de informação: \textbf{4.754}.
                \item A maior fatia é \textbf{suporte e manutenção}, não
                      desenvolvimento. Isso muda a oferta que faz sentido.
            \end{itemize}
        \end{column}
    \end{columns}
\end{frame}

\begin{frame}{Atividade 01 --- quantas ainda existem}
    \begin{columns}[T]
        \begin{column}{0.44\textwidth}
            \begin{table}
                \footnotesize
                \begin{tabular}{lrr}
                    \toprule
                    Situação & Empresas & \% \\
                    \midrule
                    Ativa    & 7.042 & 46,7 \\
                    Baixada  & 5.821 & 38,6 \\
                    Inapta   & 2.091 & 13,9 \\
                    Suspensa &   132 &  0,9 \\
                    Nula     &     9 &  0,1 \\
                    \midrule
                    Total    & \recorte & 100 \\
                    \bottomrule
                \end{tabular}
            \end{table}
        \end{column}
        \begin{column}{0.52\textwidth}
            \begin{itemize}
                \item \textbf{Metade do recorte não está ativa.} Ligar para a
                      lista filtrada é perder o dia.
                \item O cadastro diz que a empresa existe, não que ela
                      funciona.
                \item MEI é só 3,4\% --- eu esperava que fosse o ruído
                      principal, e não é.
                \item Por isso o sistema ordena, em vez de apenas filtrar.
            \end{itemize}
        \end{column}
    \end{columns}
\end{frame}

\section{Atividade 02 --- o pipeline}

\begin{frame}[fragile]{Atividade 02 --- o caminho do dado}
    \begin{columns}[T]
        \begin{column}{0.47\textwidth}
\begin{semiverbatim}\footnotesize
arquivos compactados da Receita
 |
 |- 1. fatiar      1 GiB por vez
 |- 2. ler         tudo como texto
 |- 3. recortar    UF=BA, CNAE 62 e 63
 |- 4. validar     5 regras
 |- 5. enriquecer  porte e MEI
 |- 6. contato     telefone, e-mail
 |- 7. ordenar     só ativas, por critério
      |
      |-> universo tratado
      |-> lista de leads
      |-> evidência da execução
\end{semiverbatim}
        \end{column}
        \begin{column}{0.49\textwidth}
            \begin{itemize}
                \item \textbf{Fatiar:} os 5 GiB compactados de
                      Estabelecimentos viram \textbf{16 GiB de texto} ao
                      descompactar. Cada pedaço é filtrado e apagado antes do
                      próximo.
                \item \textbf{Ler como texto:} preserva os zeros à esquerda do
                      CNPJ.
                \item \textbf{Contato:} montado no fim, para a lista já sair
                      acionável.
                \item Dois comandos reproduzem tudo, do download à planilha.
            \end{itemize}
        \end{column}
    \end{columns}
\end{frame}

\begin{frame}{Atividade 02 --- as cinco validações}
    \begin{table}
        \small
        \begin{tabular}{llr}
            \toprule
            Regra & Defeito que ela pega & Reprovados \\
            \midrule
            Dígito verificador do CNPJ & chave corrompida & 0 \\
            Situação cadastral conhecida & código fora da tabela & 0 \\
            Atividade existe na CNAE & código inexistente & 0 \\
            Data de abertura válida & o vazio \texttt{00000000} & 0 \\
            CNPJ sem repetição & ligação duplicada & 0 \\
            \bottomrule
        \end{tabular}
    \end{table}

    \begin{itemize}
        \item Nenhuma linha caiu aqui: a base é bem formada.
        \item O corte de verdade veio da situação cadastral --- e é ela que
              define quem entra na lista.
    \end{itemize}
\end{frame}

\begin{frame}{Atividade 02 --- como a lista é ordenada}
    \begin{columns}[T]
        \begin{column}{0.50\textwidth}
            \begin{table}
                \footnotesize
                \begin{tabular}{rl}
                    \toprule
                    Ordem & Critério \\
                    \midrule
                    1 & porte acima de microempresa \\
                    2 & empresa de pequeno porte \\
                    3 & não é MEI \\
                    4 & tecnologia é a atividade-fim \\
                    5 & é matriz, não filial \\
                    6 & mais de 3 anos de atividade \\
                    7 & tem nome fantasia \\
                    \bottomrule
                \end{tabular}
            \end{table}
            \footnotesize
            Empate remanescente: a mais antiga primeiro.
        \end{column}
        \begin{column}{0.46\textwidth}
            \begin{itemize}
                \item \textbf{O primeiro critério manda}; os seguintes só
                      desempatam.
                \item Estar ativa não ordena: é condição de entrada.
                \item Mudar a prioridade é mudar a ordem desta lista --- uma
                      linha, à vista de todos.
            \end{itemize}
        \end{column}
    \end{columns}
\end{frame}

\begin{frame}{Atividade 02 --- a justificativa de cada lead}
    \begin{itemize}
        \item O que ordena é o que aparece escrito: posição e motivo não
              divergem.
        \item É a última coluna da planilha, preenchida em todas as \leads{}
              linhas.
    \end{itemize}

    \begin{block}{Primeira linha da planilha}
        \footnotesize
        ABA --- COMPUTADORES, SISTEMAS E SERVICOS LTDA --- Salvador ---
        telefone e e-mail \\[1mm]
        porte acima de microempresa $|$ não é MEI $|$ tecnologia como
        atividade-fim $|$ matriz $|$ mais de 3 anos $|$ nome fantasia
    \end{block}
\end{frame}

\begin{frame}{O que o pipeline entrega}
    \begin{table}
        \small
        \begin{tabular}{llr}
            \toprule
            Arquivo & Conteúdo & Linhas $\times$ colunas \\
            \midrule
            \texttt{universo\_tratado.parquet} & recorte validado e enriquecido & 15.095 $\times$ 33 \\
            \texttt{leads\_priorizados.csv}    & a lista do analista            &  7.042 $\times$ 16 \\
            \texttt{evidencia.json}            & contagens de cada etapa        & --- \\
            \bottomrule
        \end{tabular}
    \end{table}

    \begin{itemize}
        \item \textbf{6.968 leads com telefone} (98,9\%) e 6.904 com e-mail.
        \item Competência 2026-09 processada em 2,5 minutos --- 8 na primeira
              vez, quando nada está em cache.
    \end{itemize}
\end{frame}

\begin{frame}{Decisões e próximos passos}
    \begin{columns}[T]
        \begin{column}{0.48\textwidth}
            \textbf{Decidi}
            \begin{itemize}
                \item Ativa virou condição de entrada.
                \item Contato entra na planilha; nada vai ao repositório.
                \item Prioridade por ordem declarada de critérios.
            \end{itemize}
        \end{column}
        \begin{column}{0.48\textwidth}
            \textbf{Próximo}
            \begin{itemize}
                \item Medir com o analista os 20 primeiros.
                \item Rever a ordem dos critérios pela oferta.
                \item Guardar em banco e expor consulta.
            \end{itemize}
        \end{column}
    \end{columns}
\end{frame}

\end{document}
